\chapter{Dataset}

A qualidade de um sistema de classificação de imagens depende diretamente do dataset utilizado no treinamento. No contexto deste projeto, a base de dados será obtida a partir de um dataset público, com destaque para repositórios acadêmicos e fontes amplamente utilizadas em pesquisa, como Kaggle e PlantVillage \cite{kaggle,plantvillage}. A escolha por um dataset conhecido facilita comparabilidade, reprodutibilidade e validação dos resultados.

\section{Origem e composição}

O dataset costuma ser organizado em classes representando diferentes doenças, além da classe saudável. Em muitos cenários, cada diretório corresponde a uma categoria, o que simplifica o carregamento supervisionado. A documentação do projeto prevê a divisão em conjuntos de treino, validação e teste, conforme prática padrão em aprendizado supervisionado.

\section{Estrutura das pastas}

A estrutura esperada para o treinamento geralmente segue o padrão:
\begin{itemize}
    \item \textit{train}: imagens usadas no ajuste dos pesos do modelo;
    \item \textit{validation}: imagens usadas para acompanhamento do desempenho durante o treinamento;
    \item \textit{test}: imagens mantidas para avaliação final.
\end{itemize}

Essa separação evita vazamento de informação e permite estimar de forma mais confiável a capacidade de generalização do modelo.

\section{Volume aproximado}

A quantidade de imagens pode variar de acordo com a classe e o dataset escolhido. Em projetos de classificação foliar, é comum trabalhar com milhares ou dezenas de milhares de imagens distribuídas entre várias classes. A dimensão exata do dataset final deve ser documentada quando a etapa de treinamento for executada.

\section{Treinamento supervisionado}

O treinamento supervisionado depende de rótulos consistentes e de uma política clara de balanceamento de classes. Em aplicações agrícolas, a variabilidade de iluminação, fundo, perspectiva e estágio de desenvolvimento da folha pode exigir técnicas adicionais de \textit{data augmentation}. A futura etapa de treinamento deverá registrar essas escolhas para manter a reprodutibilidade científica.

\begin{table}[H]
    \centering
    \caption{Estrutura esperada do dataset}
    \label{tab:dataset-estrutura}
    \begin{tabular}{p{3cm}p{9cm}}
        \toprule
        \textbf{Conjunto} & \textbf{Finalidade} \\
        \midrule
        Train & Ajuste dos parâmetros do modelo \\
        Validation & Monitoramento de generalização durante o treinamento \\
        Test & Avaliação final independente \\
        \bottomrule
    \end{tabular}
\end{table}
